\documentclass[10pt,a4paper]{amsart}
\usepackage[margin=19mm]{geometry}
\usepackage{amsmath,amssymb,mathtools}
\usepackage[scr=rsfs,cal=euler]{mathalfa}
\usepackage{xfrac}
\usepackage[unicode,hidelinks,bookmarks=false]{hyperref}
\newcommand{\ie}{i.e.\ }
\newcommand{\cf}{cf.\ }
\newcommand{\resp}{respectively\ }
\newcommand{\Subsection}{\subsection}
\providecommand{\nobreakdash}{\nobreak\hbox{-}}
\setlength{\emergencystretch}{2em}
\multlinegap0pt
\usepackage{bm}
\usepackage{bbm}
\usepackage{dsfont}
\usepackage{xspace}
\usepackage{cancel}
\usepackage{mathtools}
\usepackage{cleveref}
\usepackage{amsthm}
\makeatletter
\@ifundefined{lemma}{\newtheorem{lemma}{Lemma}}{}
\@ifundefined{theorem}{\newtheorem{theorem}{Theorem}}{}
\makeatother
\title[High-order Adaptive Rank Integrators for Multi-scale Linear ]{High-order Adaptive Rank Integrators for Multi-scale Linear Kinetic Transport Equations in the Hierarchical Tucker Format}
\author{William A. Sands, Wei Guo, Jing-Mei Qiu, Tao Xiong}
\date{Source: arXiv:2406.19479v3 (CC BY 4.0)}
\begin{document}
\maketitle

\noindent\emph{Source and licence.} High-order Adaptive Rank Integrators for Multi-scale Linear Kinetic Transport Equations in the Hierarchical Tucker Format. Version arXiv:2406.19479v3 (CC BY 4.0), distributed under the Creative Commons Attribution 4.0 International licence, as recorded in the arXiv licence metadata for this exact version and shown on the abstract page https://arxiv.org/abs/2406.19479v3. Full rights notice: licenses/arxiv-2406.19479-CC-BY-4.0.txt.\par\smallskip

\noindent\textit{Editorial scope.} Counted unit: the Lemma whose source label is \texttt{lem:unsplit asymptotic error bound} in the article (source0.tex lines 647--654, source label \texttt{lem:unsplit asymptotic error bound}), delivered together with the article's own complete original proof of it (source0.tex lines 655--680, one \texttt{proof} environment of 26 lines). No mathematics is written, repaired or re-derived: statement and proof are the article's own text line for line. Required context retained uncounted: eq:macro (lines 322--325); eq:micro (lines 322--325); eq:diffusion limit (lines 347--350); eq:IMEX rho final (lines 408--412); eq: IMEX g stage (lines 414--418); eq: IMEX rho stage (lines 414--418); eq:IMEX g stage explicit (lines 420--423); thm:MM AP property IMEX (lines 427--430). Every definition, display and label of those lines is retained unchanged. Omitted from this package and not redistributed: 1--321, 326--346, 351--407, 413--413, 419--419, 424--426, 431--646, 681--1417. Nothing inside a retained excerpt is omitted or altered: every retained body is delivered exactly as the article's lines are written, so no omission receipt of any kind accompanies this package. Citations: the retained excerpts carry no citation command at all, so no bibliography entry of the article is transcribed and no reference is redistributed. Reference adaptation: none was needed and none was made; every reference inside the delivered unit resolves inside it, so no unresolved reference or citation is produced. Printed slips kept literal: no printed word, symbol, hypothesis or number of the counted statement or of its proof was altered. Layout and class: the article's own class is \texttt{siamart220329} with packages including lipsum, amsfonts, amssymb, graphicx, epstopdf, algorithmic, enumitem, array, subfig, tikz, calc; the wrapper is the lane's pinned \texttt{amsart} editorial wrapper at 10pt on A4, which restates the theorem-like environments the retained text needs and adds the article's own macro definitions that the retained text needs (none), with extra wrapper packages (bm, bbm, dsfont, xspace, cancel, mathtools, cleveref). The delivered numbering is the unit's own. Assets: no retained span contains a figure environment, a graphics-inclusion command, a PDF-page inclusion, a table or a TikZ picture; no image file is copied into this package, the package declares no asset dependency and no image file is redistributed with it. Licence: version arXiv:2406.19479v3 of this article is distributed under the Creative Commons Attribution 4.0 International licence, as recorded in the arXiv licence metadata for this exact version and shown on the abstract page https://arxiv.org/abs/2406.19479v3. A full rights notice is delivered at licenses/arxiv-2406.19479-CC-BY-4.0.txt. Characters: every retained excerpt is pure ASCII, so the delivered engine, XeLaTeX with its default Latin Modern text font, reports no missing character.
\begin{align}
    &\partial_{t} \rho + \nabla_{x} \cdot \left( \left\langle \mathbf{\Omega} g \right\rangle_{\mathbf{\Omega}} \right) = - \sigma_{a} \rho + Q, \label{eq:macro} \\
    &\partial_{t} g + \frac{1}{\epsilon} \Big( I - \Pi \Big) \Big( \nabla_{x} \cdot \left(  \mathbf{\Omega} g \right) \Big) + \frac{1}{\epsilon^{2}}\nabla_{x} \cdot \left( \mathbf{\Omega} \rho \right) = -\frac{\sigma_{s}}{\epsilon^{2}} g - \sigma_{a} g, \label{eq:micro}
\end{align}

    \begin{equation}
        \label{eq:diffusion limit}
        \partial_t \rho - \nabla_x \cdot \left( \frac{1}{3 \sigma_s} \nabla_x \rho \right) = - \sigma_a \rho + Q.
    \end{equation}

\begin{align}
    \rho^{n+1} &= \rho^{n}  + \Delta t \sum_{\ell=1}^{s} \widetilde{b}_{\ell} \Bigg( -\nabla_{x} \cdot \left( \left\langle \mathbf{\Omega} g^{(\ell)} \right\rangle_{\mathbf{\Omega}} \right) - \sigma_{a} \rho^{(\ell)} + Q^{(\ell)} \Bigg), \label{eq:IMEX rho final} \\
    g^{n+1} &= g^{n} - \frac{\Delta t}{\epsilon} \sum_{\ell=1}^{s} \widetilde{b}_{\ell} \Big( I - \Pi \Big) \Big( \nabla_{x} \cdot \left(  \mathbf{\Omega} g^{(\ell)} \right) \Big) \label{eq:IMEX g final} \\ 
    &\qquad - \Delta t \sum_{\ell=1}^{s} \widetilde{b}_{\ell} \sigma_{a} g^{(\ell)} - \frac{\Delta t}{\epsilon^{2}} \sum_{\ell=1}^{s} b_{\ell} \Big( \nabla_{x} \cdot \left( \mathbf{\Omega} \rho^{(\ell)} \right) + \sigma_{s} g^{(\ell)} \Big). \nonumber
\end{align}

\begin{align}
    \rho^{(\ell)} &= \rho^{n}  + \Delta t \sum_{j=1}^{\ell-1} \widetilde{a}_{\ell j} \Big( -\nabla_{x} \cdot \left( \left\langle \mathbf{\Omega} g^{(j)} \right\rangle_{\mathbf{\Omega}} \right) - \sigma_{a} \rho^{(j)} + Q^{(j)} \Big), \label{eq: IMEX rho stage} \\
    g^{(\ell)} &= g^{n} - \frac{\Delta t}{\epsilon} \sum_{j=1}^{\ell-1} \widetilde{a}_{\ell j} \Big( I - \Pi \Big) \Big( \nabla_{x} \cdot \left(  \mathbf{\Omega} g^{(j)} \right) \Big) \label{eq: IMEX g stage} \\ 
    &\qquad - \Delta t \sum_{j=1}^{\ell-1} \widetilde{a}_{\ell j} \sigma_{a} g^{(j)} - \frac{\Delta t}{\epsilon^2} \sum_{j=1}^{\ell} a_{\ell j} \Bigg( \nabla_{x} \cdot \left( \mathbf{\Omega} \rho^{(j)} \right) + \sigma_{s} g^{(j)} \Bigg). \nonumber
\end{align}

\begin{multline}
    \label{eq:IMEX g stage explicit}
    g^{(\ell)} = \frac{1}{\epsilon^{2} + a_{\ell \ell} \Delta t \sigma_{s}} \Bigg[ \epsilon^{2} g^{n} - \epsilon \Delta t \sum_{j=1}^{\ell-1} \widetilde{a}_{\ell j} \Big( I - \Pi \Big) \Big( \nabla_{x} \cdot \left(  \mathbf{\Omega} g^{(j)} \right) \Big) - \epsilon^{2} \Delta t \sum_{j=1}^{\ell-1} \widetilde{a}_{\ell j} \sigma_{a} g^{(j)} \\  - \Delta t \sum_{j=1}^{\ell} a_{\ell j} \nabla_{x} \cdot \left( \mathbf{\Omega} \rho^{(j)} \right) -  \Delta t \sum_{j=1}^{\ell-1} a_{\ell j} \sigma_{s} g^{(j)} \Bigg].
\end{multline}

\begin{theorem}
    \label{thm:MM AP property IMEX}
    If the discretizations for the spatial and angular variables are exact and the initial data is well-prepared, then the GSA IMEX discretization of the macro-micro decomposition, given by equations \eqref{eq:IMEX rho final}-\eqref{eq: IMEX g stage}, reduces to a consistent semi-discrete approximation of the linear diffusion equation \eqref{eq:diffusion limit} as $\epsilon \rightarrow 0$.
\end{theorem}

\begin{lemma}
\label{lem:unsplit asymptotic error bound}
    Let $\boldsymbol{\rho}^{(\ell)}$ and $\mathbf{g}^{(\ell)}$ denote the solution tensors at a given stage $\ell = 1,2, \cdots s$ of the system \eqref{eq:macro}-\eqref{eq:micro} that is discretized using a GSA IMEX method. Also, let $\widetilde{\mathbf{g}}^{(\ell)}$ denote the low-rank decomposition of $\mathbf{g}^{(\ell)}$ in the HTT format with an unsplit dimension tree $\mathcal{T}$. Suppose, in addition, that no truncation is applied to the tensor $\boldsymbol{\rho}^{(\ell)}$ and that the initial data are well-prepared with $\nabla_{x}g$ remaining bounded. If the spatial and angular discretizations are exact, then the solution at each stage satisfies the inequality
    $$
        \left\lVert -\frac{1}{3\sigma_s} \nabla_x \boldsymbol{\rho}^{(\ell)} - \left\langle \mathbf{\Omega} \widetilde{\mathbf{g}}^{(\ell)} \right\rangle_{\mathbf{\Omega}} \right\rVert \lesssim \epsilon + (2 + \epsilon)\max_{\substack{1 \leq j < \ell \\ t \in \mathcal{T}}} \sigma_{t,r_{t}+1}^{(j)},
    $$
    where $\sigma_{t,r_{t}+1}^{(j)}$ denotes the leading singular value of the mode $t$ matricization of $\mathbf{g}$ at stage $j$ which is removed by truncation.
\end{lemma}

\begin{proof}
    Using the triangle inequality, we can write
    \begin{align*}
        \Bigg\lVert -\frac{1}{3\sigma_s} \nabla_x \boldsymbol{\rho}^{(\ell)} - \left\langle \mathbf{\Omega} \widetilde{\mathbf{g}}^{(\ell)} \right\rangle_{\mathbf{\Omega}} \Bigg\rVert &\leq \Bigg\lVert -\frac{1}{3\sigma_s} \nabla_x \boldsymbol{\rho}^{(\ell)} - \left\langle \mathbf{\Omega} \mathbf{g}^{(\ell)} \right\rangle_{\mathbf{\Omega}} \Bigg\rVert \\ &+ \Bigg\lVert \left\langle \mathbf{\Omega} \mathbf{g}^{(\ell)} \right\rangle_{\mathbf{\Omega}} - \left\langle \mathbf{\Omega}\widetilde{\mathbf{g}}^{(\ell)} \right\rangle_{\mathbf{\Omega}} \Bigg\rVert.
    \end{align*}
    The first term clearly satisfies
    \begin{equation*}
        \left\lVert -\frac{1}{3\sigma_s} \nabla_x \boldsymbol{\rho}^{(\ell)} - \left\langle \mathbf{\Omega} \mathbf{g}^{(\ell)} \right\rangle_{\mathbf{\Omega}} \right\rVert \lesssim \epsilon, \quad \ell = 1,2, \cdots, s,
    \end{equation*}
    since the full-grid scheme is AP in the sense of \cref{thm:MM AP property IMEX}.

    The treatment of the second term requires a bit more care. First, consider the stage update \eqref{eq: IMEX rho stage}. Let $\widetilde{\boldsymbol{\rho}}^{(\ell)}$ denote the density obtained from the update \eqref{eq: IMEX rho stage} with $\widetilde{\mathbf{g}}^{(\ell)}$ in place of $\mathbf{g}^{(\ell)}$. Then, for each $\ell = 1, 2, \cdots, s$, the inequality
    \begin{align}
        \Big\lVert \boldsymbol{\rho}^{(\ell)} - \widetilde{\boldsymbol{\rho}}^{(\ell)} \Big\rVert &\lesssim \Big\lVert \boldsymbol{\rho}^{n} - \widetilde{\boldsymbol{\rho}}^{n} \Big\rVert  + \sum_{j=1}^{\ell-1} \Big\lVert \mathbf{g}^{(j)} - \widetilde{\mathbf{g}}^{(j)} \Big\rVert + \sum_{j=1}^{\ell-1} \Big\lVert \boldsymbol{\rho}^{(j)} - \widetilde{\boldsymbol{\rho}}^{(j)} \Big\rVert, \label{eq:rho stage bound} \\
        &\lesssim \max_{\substack{1 \leq j < \ell \\ t \in \mathcal{T}}} \sigma_{t,r_{t}+1}^{(j)}, \nonumber
    \end{align}
    holds, where $\sigma_{t,r_{t}+1}^{(j)}$ denotes the leading singular value of the mode $t$ matricization of $\mathbf{g}$ at stage $j$ which is removed by truncation. This argument relies on the boundedness of angular projections and $\nabla_{x}g$.

    Next, we apply a similar argument to the stage update \eqref{eq:IMEX g stage explicit}, and with the aid of the inequality \eqref{eq:rho stage bound}, we obtain
    \begin{align*}
        \Big\lVert \left\langle \mathbf{\Omega} \mathbf{g}^{(\ell)} \right\rangle_{\mathbf{\Omega}} - \left\langle \mathbf{\Omega}\widetilde{\mathbf{g}}^{(\ell)} \right\rangle_{\mathbf{\Omega}} \Big\rVert &\lesssim \Big\lVert \mathbf{g}^{(\ell)} - \widetilde{\mathbf{g}}^{(\ell)} \Big\rVert, \\
        & \lesssim \epsilon^{2}  \Big\lVert \mathbf{g}^{n} - \widetilde{\mathbf{g}}^{n} \Big\rVert + (1 + \epsilon) \sum_{j=1}^{\ell-1} \Big\lVert \mathbf{g}^{(j)} - \widetilde{\mathbf{g}}^{(j)} \Big\rVert \\ &+ \epsilon^{2} \sum_{j=1}^{\ell-1}  \Big\lVert \mathbf{g}^{(j)} - \widetilde{\mathbf{g}}^{(j)} \Big\rVert + \sum_{j=1}^{\ell}\Big\lVert \boldsymbol{\rho}^{(j)} - \widetilde{\boldsymbol{\rho}}^{(j)} \Big\rVert, \\
        & \lesssim (2 + \epsilon) \max_{\substack{1 \leq j < \ell \\ t \in \mathcal{T}}} \sigma_{t,r_{t}+1}^{(j)}.
    \end{align*}
    Combining inequalities together yields the conclusion, which completes the proof.
\end{proof}

\end{document}
